% Two motivating examples for degree normalization: (a) two-clique dumbbell bottleneck, (b) super-hub degree dilution.
\begin{tikzpicture}[x=1mm,y=1mm,font=\figurefont\fontsize{8}{10}\selectfont,>=Stealth,
  vertex/.style={circle,draw=FigureBlue,fill=FigureBlueFill,inner sep=0pt,
    minimum size=4.6mm,line width=.55pt},
  rvertex/.style={circle,draw=FigureGreen,fill=FigureGreenFill,inner sep=0pt,
    minimum size=4.6mm,line width=.55pt},
  hub/.style={circle,draw=FigureYellow,fill=FigureYellowFill,inner sep=0pt,
    minimum size=6.2mm,line width=.65pt},
  edge/.style={draw=FigureStroke,line width=.55pt},
  stay/.style={->,draw=FigureBlue,line width=.75pt},
  bridge/.style={->,draw=FigureRed,line width=.85pt},
  note/.style={font=\FigureNoteFont\fontsize{7}{8.5}\selectfont,text=FigureInk,inner sep=1pt},
  panel/.style={font=\figurefont\fontsize{7.5}{9}\selectfont,text=FigureInk}]
  \definecolor{FigureBlue}{HTML}{7188A3}
  \definecolor{FigureBlueFill}{HTML}{E5EAF0}
  \definecolor{FigureGreen}{HTML}{829C8B}
  \definecolor{FigureGreenFill}{HTML}{E8EEE9}
  \definecolor{FigureYellow}{HTML}{B4A16F}
  \definecolor{FigureYellowFill}{HTML}{F1ECDD}
  \definecolor{FigureRed}{HTML}{A7787D}
  \definecolor{FigureInk}{HTML}{3B4252}
  \definecolor{FigureMuted}{HTML}{6B7280}
  \definecolor{FigureStroke}{HTML}{4C4D4F}

  % Panel (a): Two-clique dumbbell graph (K_m -- K_m)
  \begin{scope}[xshift=0mm]
    \fill[FigureBlueFill!55,rounded corners=2.2pt] (-4,-1) rectangle (18,21);
    \fill[FigureGreenFill!55,rounded corners=2.2pt] (28,-1) rectangle (50,21);
    \node[vertex] (a1) at (0,17) {};
    \node[vertex] (a2) at (0,3) {};
    \node[vertex] (a3) at (8,18) {};
    \node[vertex] (u)  at (14,10) {$u$};

    \node[rvertex] (v)  at (32,10) {$v$};
    \node[rvertex] (b3) at (38,18) {};
    \node[rvertex] (b1) at (46,17) {};
    \node[rvertex] (b2) at (46,3) {};

    % Clique internal edges (non-u)
    \draw[edge,draw=FigureMuted] (a1)--(a2) (a1)--(a3) (a2)--(a3);
    \draw[edge,draw=FigureMuted] (b1)--(b2) (b1)--(b3) (b2)--(b3) (v)--(b1) (v)--(b2) (v)--(b3);

    % Edges from u inside left clique
    \draw[stay] (u)--(a1);
    \draw[stay] (u)--(a2);
    \draw[stay] (u)--(a3);

    % Bridge edge u -> v
    \draw[bridge] (u)--node[above,note,text=FigureRed] {$\tfrac{1}{m}$} (v);
    \node[note,text=FigureBlue] at (7,23.5) {退回团内 $\tfrac{m-1}{m}$};
    \node[note,text=FigureMuted] at (6,-3.5) {$K_m$};
    \node[note,text=FigureMuted] at (40,-3.5) {$K_m$};
    \node[panel] at (23,-9.5) {(a) 双团哑铃图的跨团瓶颈};
  \end{scope}

  % Panel (b): Super-hub vs low-degree vertex
  \begin{scope}[xshift=64mm]
    \node[hub] (h) at (12,10) {$h$};
    \node[vertex,minimum size=3.8mm] (p1) at (2,19) {};
    \node[vertex,minimum size=3.8mm] (p2) at (0,10) {};
    \node[vertex,minimum size=3.8mm] (p3) at (2,1) {};
    \node[vertex,minimum size=3.8mm] (p4) at (12,20) {};
    \node[vertex,minimum size=3.8mm] (p5) at (23,17) {};
    \node[vertex,minimum size=3.8mm] (p6) at (23,3) {};
    \node[rvertex] (i) at (36,10) {$i$};

    \draw[stay] (h)--(p1);
    \draw[stay] (h)--(p2);
    \draw[stay] (h)--(p3);
    \draw[stay] (h)--(p4);
    \draw[stay] (h)--node[above,note,text=FigureBlue] {$\tfrac{1}{D}$} (p5);
    \draw[stay] (h)--node[below,note,text=FigureBlue] {$\tfrac{1}{D}$} (p6);

    \draw[bridge,draw=FigureGreen] (i)--node[above right,note,text=FigureGreen] {$\tfrac{1}{2}$} (p5);
    \draw[bridge,draw=FigureGreen] (i)--node[below right,note,text=FigureGreen] {$\tfrac{1}{2}$} (p6);

    \node[note,text=FigureInk] at (12,-3.5) {度数 $d_h=D\gg 1$};
    \node[note,text=FigureInk] at (36,-3.5) {度数 $d_i=2$};
    \node[panel] at (19,-9.5) {(b) 超级节点稀释单边转移概率};
  \end{scope}
\end{tikzpicture}
